\documentclass[11pt]{article}
\usepackage[margin=1in]{geometry}
\usepackage{hyperref}
\usepackage{listings}
\usepackage{xcolor}
\usepackage{amsmath}

\lstset{
  basicstyle=\ttfamily\small,
  breaklines=true,
  columns=fullflexible,
  frame=single,
  backgroundcolor=\color{gray!10},
  showstringspaces=false
}

\title{URI OPeNDAP Test Case: Working Log (Prompt \#5)}
\author{Compiled for P. Cornillon}
\date{2026-05-16}

\begin{document}
\maketitle

\section*{Purpose}

Working log for Prompt~\#5: write user-facing descriptions of the
two readable datasets, plus the two agent-facing usage files.

\section{Prompt \#5 (2026-05-16 22:06 UTC)}

\subsection{Instruction}

\emph{``Write a description of each of the two datasets for a
typical user; i.e., not too detailed but detailed enough for the
user to decide whether or not the dataset meets their needs and to
acquire subsets of the data in which they might be interested.
Also, provide \texttt{usage\_sst.md} and
\texttt{usage\_gradient\_SST.md} files, which the user can provide
to an AI agent to acquire and make use of the data.''}

\subsection{Deliverables}

Four files (besides this log):

\begin{itemize}
  \item \texttt{docs/uri\_test\_case\_5.tex} --- short, accessible
    descriptions of \texttt{SST\_Orbits/} and
    \texttt{gradients\_by\_period/}.
  \item \texttt{docs/usage\_sst.md} --- self-contained markdown
    guide for the orbit branch, intended to be handed to an AI
    agent or a working scientist.
  \item \texttt{docs/usage\_gradient\_SST.md} --- the analog for
    the gradient-statistics branch, with explicit notes on every
    inferred attribute.
  \item \texttt{scripts/usage\_sanity.py} --- sanity test that
    runs the headline code examples from both \texttt{.md} files
    against the live server, to verify they work as written.
\end{itemize}

\subsection{Drafting choices}

\paragraph{The user description (uri\_test\_case\_5.tex).}
Kept short. For each dataset, a six-line \emph{at-a-glance} block,
a \emph{who it is for} paragraph, a $\sim$15-line code snippet for
a sub-set, and an \emph{honest limitations} bullet list. The
deeper material (full variable tables, exhaustive caveats,
algorithm-support estimate) is in earlier outputs of this test
case and in the two \texttt{.md} files.

\paragraph{usage\_sst.md.} Structured for a Python-fluent agent.
Server URL, filename pattern, DAP4-preferred access via
\texttt{pydap.client.open\_url}, variable inventory grouped by
role (SST/gradients, geolocation, time, AMSR-E co-located,
provenance), code recipes for the most common cases (read a
window, get the orbit time, bracket by nadir lat), and a numbered
list of the metadata pitfalls a careful agent must handle
(non-UDUNITS units, scale\_factor on Int32 lat/lon, missing time
reference in \texttt{DateTime}, malformed \texttt{flag\_masks},
AMSR-E placeholders after 2011, the two clearly wrong unit
declarations in the L2eqa group).

\paragraph{usage\_gradient\_SST.md.} Structured the same way but
with everything declared up-front as \emph{inferred} because the
file itself carries no metadata. Includes the verified
cell-centered grid (lat\_idx 0 = $-89.5^\circ$, lon\_idx 0 =
$-179.5^\circ$), the per-variable units (all derived from the
naming grammar and the sibling \texttt{SST\_Orbits/} branch), the
divisor rule (\texttt{*\_sum\_grad\_as\_per\_km} pairs with
\texttt{as\_pixel\_count}, \texttt{*\_sum\_grad\_at\_per\_km} with
\texttt{at\_pixel\_count}, \texttt{*\_sum\_grad\_mag\_per\_km}
with \texttt{at\_pixel\_count} verified empirically), and a
\S6.3 ``what I won't guess'' section that lists items where the
agent must defer rather than fabricate.

\subsection{Sanity check}

\texttt{scripts/usage\_sanity.py} ran the three headline examples
from the two markdown files against the live server. Outputs:

\begin{lstlisting}
SST_In sample (orbit 115220, eq, j=600:700)
  n_valid=10000, range=[16.16,29.19] C, mean=27.53 C
DateTime -> 2024-01-01T01:15:58 UTC

gradients_by_period 07_2020:
  per-cell day-mean SST     n=38540, range=[-1.12,35.17], mean=16.75
  per-cell day-mean |grad|  n=39909, range=[0.001,0.84], mean=0.046

Bounding box (Atlantic 100W..50W, 20N..50N):
  shape (51,31), n_valid=1032, mean SST=25.28 C
\end{lstlisting}

All three return sensible oceanographic values. One example in
\texttt{usage\_sst.md} originally read SST at line index
20000:20100 of orbit 115220; that index range turned out to be in
the polar Kara Sea region (every pixel below the
\texttt{valid\_min}, returning empty after masking), so the
example was rewritten to read at line indices 29200:29300 (near
the equator) and a tip was added telling the reader how to pick a
window with valid SST by bracketing on \texttt{nadir\_latitude}.

The second small fix: the \texttt{pydap} slice object returned by
\texttt{ds["SST\_In"][...]} is a \texttt{SelfClearingArray}, not a
numpy ndarray. The example now wraps the slice in
\texttt{np.asarray(\ldots .data, dtype="f8")} before the
\texttt{.astype("f8")} step. (This was caught by the sanity
script.)

\subsection{Files produced}

\begin{lstlisting}
docs/uri_test_case_5.tex
docs/test_case_uri_log_5.tex      (this file)
docs/usage_sst.md
docs/usage_gradient_SST.md
scripts/usage_sanity.py
\end{lstlisting}

\subsection{Session timing}

Started 2026-05-16 22:06\,UTC. Closed log 2026-05-16
$\sim$22:50\,UTC.

\end{document}
